\documentclass{ximera}

\input{../../../preamble.tex}


\definecolor{fvdnavy}{HTML}{071D43}
\definecolor{fvdcyan}{HTML}{20BFC6}
\definecolor{fvdcyanlight}{HTML}{EFFCFB}
\definecolor{fvdmagenta}{HTML}{F10D69}
\definecolor{fvdmagentalight}{HTML}{FFF1F7}
\definecolor{fvdtext}{HTML}{163044}





%=================================================
% Pedagogische omgevingen
% PDF: tcolorbox
% HTML: learning-box
%=================================================

\newenvironment{voorkennis}
{%
    \ifdefined\HCode
        \HCode{%
<section class="learning-box learning-box--prior">
<div class="learning-box__header">
<span class="learning-box__icon" aria-hidden="true"></span>
<span class="learning-box__title">Voorkennis</span>
</div>
<div class="learning-box__content">
        }%
        \begin{itemize}
    \else
       \begin{tcolorbox}[
    enhanced,
    breakable,
    colback=fvdcyanlight,
    colframe=fvdcyan,
    coltext=fvdtext,
    boxrule=0.5pt,
    leftrule=4pt,
    arc=4mm,
    outer arc=4mm,
    left=7mm,
    right=6mm,
    top=5mm,
    bottom=5mm,
    before skip=6mm,
    after skip=6mm,
    title=Voorkennis,
    coltitle=fvdcyan,
    fonttitle=\bfseries\Large,
    attach title to upper={\par\smallskip},
    boxed title style={
        empty,
        left=0mm,
        right=0mm,
        top=0mm,
        bottom=0mm
    },
    overlay={
        \draw[
            fvdcyan,
            line width=0.5pt
        ]
        ([xshift=42mm,yshift=-13mm]frame.north west)
        --
        ([xshift=-7mm,yshift=-13mm]frame.north east);
    }
]
        \begin{itemize}
    \fi
}
{%
    \end{itemize}
    \ifdefined\HCode
        \HCode{%
</div>
</section>
        }%
    \else
        \end{tcolorbox}
    \fi
}


\newenvironment{leerdoelen}
{%
    \ifdefined\HCode
        \HCode{%
<section class="learning-box learning-box--goals">
<div class="learning-box__header">
<span class="learning-box__icon" aria-hidden="true"></span>
<span class="learning-box__title">Leerdoelen</span>
</div>
<div class="learning-box__content">
        }%
        \begin{itemize}
    \else
\begin{tcolorbox}[
    enhanced,
    breakable,
    colback=fvdmagentalight,
    colframe=fvdmagenta,
    coltext=fvdtext,
    boxrule=0.5pt,
    leftrule=4pt,
    arc=4mm,
    outer arc=4mm,
    left=7mm,
    right=6mm,
    top=5mm,
    bottom=5mm,
    before skip=6mm,
    after skip=6mm,
    title=Leerdoelen,
    coltitle=fvdmagenta,
    fonttitle=\bfseries\Large,
    attach title to upper={\par\smallskip},
    boxed title style={
        empty,
        left=0mm,
        right=0mm,
        top=0mm,
        bottom=0mm
    },
    overlay={
        \draw[
            fvdmagenta,
            line width=0.5pt
        ]
        ([xshift=42mm,yshift=-13mm]frame.north west)
        --
        ([xshift=-7mm,yshift=-13mm]frame.north east);
    }
]
        \begin{itemize}
    \fi
}
{%
    \end{itemize}
    \ifdefined\HCode
        \HCode{%
</div>
</section>
        }%
    \else
        \end{tcolorbox}
    \fi
}


\addPrintStyle{../../..}

\title{Bewegen in twee dimensies}
\author{Ingmar Herreman}

\begin{abstract}
Tot nu toe beschreven we bewegingen langs één rechte. Veel bewegingen
in de werkelijkheid zijn echter tweedimensionaal. Een bal die
horizontaal van een tafel wordt weggeschoten, beweegt tegelijk
horizontaal én verticaal.

In dit hoofdstuk bouwen we daarom de vectoriële beschrijving van
beweging verder uit. We beschrijven positie, verplaatsing, snelheid en
versnelling met componenten en onderzoeken hoe twee onafhankelijke
componentbewegingen samen één baan vormen.

De horizontale worp is onze centrale toepassing. Horizontaal voert het
voorwerp een eenparige rechtlijnige beweging uit; verticaal een vrije
val. Door beide bewegingen samen te brengen kunnen we positie,
snelheid, valtijd, dracht en baan analyseren en kwantificeren.

We onderzoeken de horizontale worp ook experimenteel en leren kritisch
omgaan met het model zonder luchtweerstand.

Voor wie verder wil gaan, eindigt het hoofdstuk met een verdieping over
de schuine worp. Daar blijkt dat de horizontale worp slechts een
bijzonder geval is van een veel algemener model voor projectielbeweging.
\end{abstract}

\begin{document}

% FVD_CHAPTER_DATA_START
% Automatisch gegenereerd uit index.tex.
% Niet handmatig aanpassen.
\fvdchapterdata{1}{Beweging en kracht — Van beweging beschrijven tot energie begrijpen}{4}
% FVD_CHAPTER_DATA_END



\maketitle

% ============================================================
% VOORKENNIS EN LEERDOELEN
% ============================================================

\begin{voorkennis}
  \item Je kan een beweging in één dimensie beschrijven met positie,
        snelheid en versnelling.
  \item Je kan een ERB en een EVRB analyseren.
  \item Je kan vrije val beschrijven met $y(t)$, $v_y(t)$ en $a_y(t)$.
  \item Je kan vectoren tekenen en ontbinden in componenten.
  \item Je kan de grootte van een vector uit zijn loodrechte componenten
        bepalen met de stelling van Pythagoras.
  \item Je kan een referentiestelsel kiezen en consequent met tekens werken.
\end{voorkennis}

\begin{leerdoelen}
  \item Je kan positie, verplaatsing, snelheid en versnelling bij een
        tweedimensionale beweging vectorieel beschrijven.
  \item Je kan vectoren correct tekenen en hun $x$- en $y$-componenten
        interpreteren.
  \item Je kan het onafhankelijkheidsbeginsel uitleggen en toepassen.
  \item Je kan een horizontale worp ontbinden in een horizontale ERB en
        een verticale vrije val.
  \item Je kan $x(t)$, $y(t)$, $v_x(t)$, $v_y(t)$, $a_x(t)$ en $a_y(t)$
        bij een horizontale worp opstellen en interpreteren.
  \item Je kan de snelheidsvector tijdens een horizontale worp construeren
        en zijn grootte en richting bepalen.
  \item Je kan valtijd, horizontale dracht en impactgrootheden bepalen.
  \item Je kan de baanvergelijking van een horizontale worp afleiden en
        interpreteren.
  \item Je kan problemen over horizontale worpen systematisch mathematiseren.
  \item Je kan een horizontale worp experimenteel onderzoeken en meetgegevens
        vergelijken met het theoretische model.
  \item Je kan de beperkingen van het ideale projectielmodel aangeven.
\end{leerdoelen}


% ============================================================
% 1 OPENING
% ============================================================

\FVDsection{Twee ballen, één verrassende vraag}

Stel je twee identieke ballen voor op dezelfde hoogte.

Bal A wordt zonder beginsnelheid losgelaten.

Bal B wordt op exact hetzelfde ogenblik horizontaal weggeschoten.

Welke bal raakt eerst de grond?

De horizontaal weggeschoten bal legt tijdens zijn val een veel langere
baan af. Het is daarom verleidelijk te denken dat hij langer onderweg
is. Toch is dat in het ideale model niet zo.

De sleutel is dat we niet naar de volledige beweging in één keer
kijken, maar naar twee loodrechte componenten.

\begin{oefening}[Voorspel vóór je rekent]{\oefB\ESS}

Twee identieke ballen vertrekken op hetzelfde ogenblik van dezelfde
hoogte. Bal A wordt losgelaten; bal B krijgt een horizontale
beginsnelheid.

Verwaarloos luchtweerstand.

\begin{question}
Welke bal verwacht je eerst op de grond? Noteer je voorspelling.
\end{question}

\begin{question}
Heeft bal B aanvankelijk een verticale snelheidscomponent?
\end{question}

\begin{question}
Vergelijk de verticale versnellingen van beide ballen.
\end{question}

Bewaar je eerste antwoord. We komen erop terug nadat we het model
hebben opgebouwd.

\end{oefening}


% ============================================================
% 2 VAN 1D NAAR 2D
% ============================================================

\FVDsection{Van één naar twee dimensies}

Bij een rechtlijnige beweging volstond één coördinaat, bijvoorbeeld
$x(t)$. Bij een beweging in een vlak hebben we in het algemeen twee
coördinaten nodig:

\[
x(t)
\qquad\text{en}\qquad
y(t).
\]

Samen bepalen ze op elk tijdstip de positie van het voorwerp.


\FVDsubsection{De positievector}

Ten opzichte van een gekozen oorsprong $O$ beschrijven we de positie
met de vector

\[
\boxed{
\vec r(t)=x(t)\,\vec e_x+y(t)\,\vec e_y.
}
\]

Hierbij zijn $\vec e_x$ en $\vec e_y$ eenheidsvectoren langs de
$x$- en $y$-as.

In componentnotatie kunnen we schrijven:

\[
\vec r(t)=
\begin{pmatrix}
x(t)\\
y(t)
\end{pmatrix}.
\]

De positievector wijst dus van de oorsprong naar de actuele positie
van het voorwerp.


\FVDsubsection{De verplaatsingsvector}

Tussen twee tijdstippen $t_1$ en $t_2$ is de verplaatsing:

\[
\boxed{
\Delta\vec r
=
\vec r(t_2)-\vec r(t_1).
}
\]

Dus:

\[
\Delta\vec r
=
\Delta x\,\vec e_x+\Delta y\,\vec e_y.
\]

De verplaatsingsvector verbindt de beginpositie rechtstreeks met de
eindpositie. Hij is niet noodzakelijk gelijk aan de afgelegde baan.


\begin{oefening}[Positie en verplaatsing]{\oefB\ESS}

Een drone bevindt zich eerst in

\[
P(2,0\,\mathrm{m},\,3,0\,\mathrm{m})
\]

en later in

\[
Q(8,0\,\mathrm{m},\,7,0\,\mathrm{m}).
\]

\begin{question}
Schrijf de twee positievectoren.
\end{question}

\begin{question}
Bepaal de verplaatsingsvector.
\end{question}

\begin{question}
Bepaal de grootte van de verplaatsing.
\end{question}

\begin{question}
Leg uit waarom je uit deze gegevens niet noodzakelijk de afgelegde weg
van de drone kent.
\end{question}

\end{oefening}


% ============================================================
% 3 SNELHEID EN VERSNELLING IN 2D
% ============================================================

\FVDsection{Snelheid en versnelling als vectoren}

De snelheidsvector beschrijft zowel hoe snel als in welke richting de
positie verandert.

\[
\boxed{
\vec v(t)=v_x(t)\,\vec e_x+v_y(t)\,\vec e_y.
}
\]

Voor leerlingen Wetenschappen-Wiskunde kunnen we dit ook schrijven als

\[
\boxed{
\vec v(t)=\vec r\,'(t).
}
\]

Componentgewijs:

\[
v_x(t)=x'(t),
\qquad
v_y(t)=y'(t).
\]

De grootte van de snelheid is:

\[
\boxed{
v=|\vec v|
=
\sqrt{v_x^2+v_y^2}.
}
\]

De richting van $\vec v$ is op elk punt raaklijnig aan de baan.


\FVDsubsection{De versnellingsvector}

Ook versnelling is een vector:

\[
\boxed{
\vec a(t)=a_x(t)\,\vec e_x+a_y(t)\,\vec e_y.
}
\]

Met afgeleiden:

\[
\boxed{
\vec a(t)=\vec v\,'(t)=\vec r\,''(t).
}
\]

Dus:

\[
a_x=v_x'(t),
\qquad
a_y=v_y'(t).
\]

Een voorwerp kan een constante snelheidscomponent in één richting
hebben en tegelijk versnellen in een loodrechte richting.


\begin{opmerking}

Dat de \emph{grootte} van de snelheid verandert, is niet de enige
mogelijkheid voor versnelling. Ook een verandering van richting van
de snelheidsvector betekent versnelling.

Dit idee wordt later essentieel bij de cirkelbeweging.

\end{opmerking}


\begin{oefening}[Van componenten naar vector]{\oefB\ESS}

Op een bepaald tijdstip heeft een voorwerp

\[
v_x=12,0\,\mathrm{m/s},
\qquad
v_y=-5,0\,\mathrm{m/s}.
\]

\begin{question}
Schrijf de snelheidsvector.
\end{question}

\begin{question}
Bepaal de grootte van de snelheid.
\end{question}

\begin{question}
Beschrijf de bewegingsrichting in woorden.
\end{question}

\begin{question}
Schets de componenten en de resulterende snelheidsvector.
\end{question}

\end{oefening}


% ============================================================
% 4 ONAFHANKELIJKHEIDSBEGINSEL
% ============================================================

\FVDsection{Het onafhankelijkheidsbeginsel}

Een tweedimensionale beweging kan vaak worden geanalyseerd als de
samenstelling van twee loodrechte bewegingen.

\begin{eigenschap}[Onafhankelijkheidsbeginsel]

De beweging in de $x$-richting en de beweging in de $y$-richting
kunnen afzonderlijk worden geanalyseerd. Beide componentbewegingen
verlopen gelijktijdig en hebben dezelfde tijdsvariabele $t$.

\end{eigenschap}

Dat betekent niet dat er twee verschillende voorwerpen bewegen. Het
is één beweging die wij voor de analyse ontbinden in twee componenten:

\[
\boxed{
\text{2D-beweging}
=
\text{$x$-beweging}
+
\text{$y$-beweging}.
}
\]

De tijd vormt de verbinding tussen beide beschrijvingen.


\FVDsubsection{Vectoriële samenstelling}

Als tijdens een klein tijdsinterval een voorwerp horizontaal
$\Delta x$ en verticaal $\Delta y$ verplaatst, dan is de totale
verplaatsing:

\[
\Delta\vec r
=
\Delta x\,\vec e_x+\Delta y\,\vec e_y.
\]

Hetzelfde geldt voor snelheid:

\[
\vec v
=
v_x\,\vec e_x+v_y\,\vec e_y.
\]

En voor versnelling:

\[
\vec a
=
a_x\,\vec e_x+a_y\,\vec e_y.
\]


\begin{oefening}[Eén beweging, twee componenten]{\oefU\ESS}

Een puntmassa heeft:

\[
x(t)=4,0t,
\qquad
y(t)=6,0-4,905t^2.
\]

\begin{question}
Beschrijf de beweging in de $x$-richting.
\end{question}

\begin{question}
Beschrijf de beweging in de $y$-richting.
\end{question}

\begin{question}
Bepaal $v_x(t)$ en $v_y(t)$.
\end{question}

\begin{question}
Bepaal $a_x(t)$ en $a_y(t)$.
\end{question}

\begin{question}
Welke bekende fysische beweging herken je?
\end{question}

\end{oefening}


% ============================================================
% 5 HORIZONTALE WORP
% ============================================================

\FVDsection{De horizontale worp}

Een \textbf{horizontale worp} ontstaat wanneer een voorwerp met een
horizontale beginsnelheid wordt gelanceerd en daarna, in het ideale
model, alleen door de zwaartekracht wordt beïnvloed.

We verwaarlozen luchtweerstand en nemen $g$ constant.

We kiezen:

\begin{itemize}
  \item de $x$-as horizontaal naar rechts;
  \item de $y$-as verticaal omhoog;
  \item het lanceerpunt op $(x_0,y_0)$.
\end{itemize}

De beginsnelheid is volledig horizontaal:

\[
\vec v_0=v_0\,\vec e_x.
\]

Dus:

\[
v_{x,0}=v_0,
\qquad
v_{y,0}=0.
\]

De versnelling is volledig verticaal:

\[
\boxed{
\vec a=-g\,\vec e_y.
}
\]

Dus:

\[
a_x=0,
\qquad
a_y=-g.
\]


\FVDsubsection{Horizontaal: een ERB}

Omdat

\[
a_x=0,
\]

blijft de horizontale snelheidscomponent constant:

\[
\boxed{
v_x(t)=v_0.
}
\]

De horizontale positie is:

\[
\boxed{
x(t)=x_0+v_0t.
}
\]


\FVDsubsection{Verticaal: vrije val}

Verticaal is de beginsnelheid nul:

\[
v_{y,0}=0.
\]

Daarom:

\[
\boxed{
v_y(t)=-gt
}
\]

en:

\[
\boxed{
y(t)=y_0-\frac12gt^2.
}
\]


\begin{eigenschap}[Horizontale worp als samenstelling]

\[
\boxed{
\begin{array}{ccl}
x\text{-richting} &:& \text{ERB}\\[1mm]
y\text{-richting} &:& \text{vrije val}
\end{array}
}
\]

De horizontale en verticale beweging delen alleen dezelfde tijd $t$.

\end{eigenschap}


\begin{oefening}[De basisvergelijkingen]{\oefB\ESS}

Een bal verlaat horizontaal een tafel met

\[
v_0=3,0\,\mathrm{m/s}.
\]

De tafel is $1,25\,\mathrm{m}$ hoog. Kies de grond als $y=0$ en het
punt loodrecht onder de tafelrand als $x=0$.

\begin{question}
Noteer $x_0$, $y_0$, $v_{x,0}$, $v_{y,0}$, $a_x$ en $a_y$.
\end{question}

\begin{question}
Stel $x(t)$ en $y(t)$ op.
\end{question}

\begin{question}
Stel $v_x(t)$ en $v_y(t)$ op.
\end{question}

\end{oefening}


% ============================================================
% 6 TWEE BALLEN
% ============================================================

\FVDsection{Waarom vallen beide ballen tegelijk?}

We keren terug naar de openingsvraag.

Bal A wordt losgelaten. Bal B wordt gelijktijdig horizontaal
weggeschoten. Beide vertrekken van dezelfde hoogte en hebben dezelfde
verticale beginsnelheid:

\[
v_{y,0}=0.
\]

Voor beide geldt verticaal:

\[
y(t)=y_0-\frac12gt^2.
\]

Hun verticale bewegingen zijn dus identiek.

De horizontale snelheid van bal B heeft geen invloed op de tijd die
nodig is om verticaal dezelfde afstand te vallen.

\begin{eigenschap}

In het ideale model raken twee voorwerpen die op hetzelfde tijdstip
van dezelfde hoogte vertrekken en dezelfde verticale beginsnelheid
hebben, tegelijk de grond, ook wanneer één van beide een horizontale
snelheidscomponent heeft.

\end{eigenschap}


\begin{oefening}[Herbekijk je voorspelling]{\oefU\ESS}

Keer terug naar de eerste oefening.

\begin{question}
Was je oorspronkelijke voorspelling correct?
\end{question}

\begin{question}
Bewijs met de verticale bewegingsvergelijking dat de valtijd niet van
$v_x$ afhangt.
\end{question}

\begin{question}
Leg hetzelfde resultaat uit zonder formules, met het
onafhankelijkheidsbeginsel.
\end{question}

\begin{question}
Welke aanname uit ons model zou in een werkelijk experiment kleine
verschillen kunnen veroorzaken?
\end{question}

\end{oefening}


% ============================================================
% 7 VALTIJD EN DRACHT
% ============================================================

\FVDsection{Valtijd en horizontale dracht}

Stel dat een voorwerp horizontaal wordt weggeschoten vanaf hoogte $h$
boven een horizontale grond.

Kies:

\[
y_0=h,
\qquad
y_{\mathrm{grond}}=0.
\]

Bij impact:

\[
0=h-\frac12gt^2.
\]

Daaruit:

\[
\boxed{
t_{\mathrm{val}}
=
\sqrt{\frac{2h}{g}}.
}
\]

Merk op: de horizontale beginsnelheid komt niet in deze formule voor.


\FVDsubsection{De horizontale dracht}

Tijdens de valtijd beweegt het voorwerp horizontaal met constante
snelheid $v_0$.

Als $x_0=0$:

\[
d=v_0t_{\mathrm{val}}.
\]

Dus:

\[
\boxed{
d
=
v_0\sqrt{\frac{2h}{g}}.
}
\]

Deze formule is nuttig, maar belangrijker is dat je haar kunt
reconstrueren vanuit de twee componentbewegingen.


\begin{voorbeeld}[Een tennisbal van een tafel]

Een tennisbal rolt horizontaal van een tafel van
$1,20\,\mathrm{m}$ hoog met een snelheid van
$4,50\,\mathrm{m/s}$.

De valtijd:

\[
t_{\mathrm{val}}
=
\sqrt{\frac{2\cdot1,20}{9,81}}
\approx0,495\,\mathrm{s}.
\]

De horizontale dracht:

\[
d
=
4,50\cdot0,495
\approx2,23\,\mathrm{m}.
\]

\end{voorbeeld}


\begin{oefening}[Van tafel naar vloer]{\oefB\ESS}

Een knikker verlaat horizontaal een tafel van
$0,80\,\mathrm{m}$ hoog met $v_0=2,50\,\mathrm{m/s}$.

\begin{question}
Bepaal de valtijd.
\end{question}

\begin{question}
Bepaal de horizontale dracht.
\end{question}

\begin{question}
Wat gebeurt er met de valtijd als de horizontale beginsnelheid
verdubbelt?
\end{question}

\begin{question}
Wat gebeurt er dan met de dracht?
\end{question}

\end{oefening}


\begin{oefening}[Schaalredenering]{\oefU\ESS}

Voor een horizontale worp geldt:

\[
d=v_0\sqrt{\frac{2h}{g}}.
\]

Onderzoek zonder telkens nieuwe getallen in te vullen:

\begin{question}
hoe $d$ verandert als $v_0$ verdubbelt;
\end{question}

\begin{question}
hoe $d$ verandert als $h$ verviervoudigt;
\end{question}

\begin{question}
hoe $d$ verandert als zowel $v_0$ als $h$ verdubbelen.
\end{question}

Motiveer vanuit de formule én vanuit de fysica.

\end{oefening}


% ============================================================
% 8 SNELHEIDSVECTOR
% ============================================================

\FVDsection{De snelheidsvector verandert voortdurend}

Bij de horizontale worp blijft:

\[
v_x=v_0.
\]

Maar:

\[
v_y=-gt.
\]

De totale snelheidsvector is dus:

\[
\boxed{
\vec v(t)
=
v_0\,\vec e_x-gt\,\vec e_y.
}
\]

Hoewel de horizontale component constant blijft, verandert de totale
snelheidsvector voortdurend.


\FVDsubsection{Grootte van de snelheid}

Omdat de componenten loodrecht op elkaar staan:

\[
\boxed{
|\vec v(t)|
=
\sqrt{v_0^2+(gt)^2}.
}
\]

De grootte van de snelheid neemt dus toe tijdens de val.


\FVDsubsection{Richting van de snelheid}

Noem $\theta$ de hoek die de snelheidsvector onder de horizontale
maakt. Dan:

\[
\tan\theta
=
\frac{|v_y|}{v_x}
=
\frac{gt}{v_0}.
\]

Dus:

\[
\boxed{
\theta
=
\arctan\left(\frac{gt}{v_0}\right).
}
\]

De snelheidsvector is steeds raaklijnig aan de baan.


\begin{voorbeeld}[Snelheid bij impact]

Voor de tennisbal uit het vorige voorbeeld:

\[
v_x=4,50\,\mathrm{m/s}.
\]

Na $0,495\,\mathrm{s}$:

\[
v_y
=
-9,81\cdot0,495
\approx-4,86\,\mathrm{m/s}.
\]

De grootte van de impactsnelheid:

\[
|\vec v|
=
\sqrt{4,50^2+4,86^2}
\approx6,62\,\mathrm{m/s}.
\]

\end{voorbeeld}


\begin{oefening}[Vectoren langs de baan]{\oefU\ESS}

Een voorwerp wordt horizontaal gelanceerd met

\[
v_0=8,0\,\mathrm{m/s}.
\]

\begin{question}
Bepaal $\vec v$ na $0,50\,\mathrm{s}$.
\end{question}

\begin{question}
Bepaal de grootte van de snelheid.
\end{question}

\begin{question}
Bepaal de hoek van de snelheid onder de horizontale.
\end{question}

\begin{question}
Schets op drie opeenvolgende punten van de baan de snelheidsvector.
Wat blijft gelijk en wat verandert?
\end{question}

\begin{question}
Schets op dezelfde punten de versnellingsvector.
\end{question}

\end{oefening}


% ============================================================
% 9 GRAFIEKEN
% ============================================================

\FVDsection{Zes functies vertellen één verhaal}

Een horizontale worp kunnen we volledig beschrijven met zes
componentfuncties.

Voor $x_0=0$ en lanceerhoogte $y_0=h$:

\[
x(t)=v_0t,
\qquad
v_x(t)=v_0,
\qquad
a_x(t)=0,
\]

en:

\[
y(t)=h-\frac12gt^2,
\qquad
v_y(t)=-gt,
\qquad
a_y(t)=-g.
\]

Horizontaal herkennen we dus:

\[
\text{lineaire positie}
\longrightarrow
\text{constante snelheid}
\longrightarrow
\text{nulversnelling}.
\]

Verticaal:

\[
\text{kwadratische positie}
\longrightarrow
\text{lineaire snelheid}
\longrightarrow
\text{constante versnelling}.
\]


\begin{oefening}[Van functies naar grafieken]{\oefU\ESS}

Voor een horizontale worp geldt:

\[
v_0=6,0\,\mathrm{m/s},
\qquad
h=5,0\,\mathrm{m}.
\]

Schets voor de volledige vlucht:

\[
x(t),\quad v_x(t),\quad a_x(t),
\]

en:

\[
y(t),\quad v_y(t),\quad a_y(t).
\]

Duid op alle grafieken het tijdstip van impact aan.

Beschrijf vervolgens minstens vier verbanden tussen de zes grafieken.

\end{oefening}


% ============================================================
% 10 BAANVERGELIJKING
% ============================================================

\FVDsection{Van tijdsfuncties naar de baan}

De vergelijkingen

\[
x(t)=x_0+v_0t
\]

en

\[
y(t)=y_0-\frac12gt^2
\]

geven de positie als functie van de tijd.

Soms willen we rechtstreeks weten welke $y$-waarde hoort bij een
gegeven $x$-waarde. Dan elimineren we de tijd.

Uit:

\[
x=x_0+v_0t
\]

volgt:

\[
t=\frac{x-x_0}{v_0}.
\]

Invullen in de verticale plaatsfunctie geeft:

\[
y
=
y_0
-
\frac12g
\left(
\frac{x-x_0}{v_0}
\right)^2.
\]

Dus:

\[
\boxed{
y(x)
=
y_0
-
\frac{g}{2v_0^2}(x-x_0)^2.
}
\]

Dit is een kwadratische functie. De ideale baan van een horizontaal
geworpen puntmassa is dus een deel van een parabool.


\begin{opmerking}

De baanvergelijking bevat geen tijd meer. Ze vertelt \emph{waar} de
baan ligt, maar niet rechtstreeks \emph{wanneer} het voorwerp een
bepaald punt bereikt.

\end{opmerking}


\begin{oefening}[Van parametervergelijkingen naar parabool]{\oefU\ESS}

Een bal wordt horizontaal gelanceerd vanuit

\[
(x_0,y_0)=(0,\,10,0\,\mathrm{m})
\]

met

\[
v_0=7,0\,\mathrm{m/s}.
\]

\begin{question}
Stel $x(t)$ en $y(t)$ op.
\end{question}

\begin{question}
Elimineer $t$ en bepaal $y(x)$.
\end{question}

\begin{question}
Bepaal met de baanvergelijking de hoogte wanneer $x=5,0\,\mathrm{m}$.
\end{question}

\begin{question}
Controleer hetzelfde punt via de tijdsfuncties.
\end{question}

\end{oefening}


\begin{oefening}[Wat verandert de parabool?]{\oefG}

Beschouw horizontale worpen vanaf hetzelfde punt met verschillende
horizontale beginsnelheden.

Gebruik

\[
y(x)
=
y_0-\frac{g}{2v_0^2}x^2.
\]

\begin{question}
Onderzoek wat er met de vorm van de baan gebeurt wanneer $v_0$
toeneemt.
\end{question}

\begin{question}
Leg dit resultaat fysisch uit.
\end{question}

\begin{question}
Gebruik eventueel ICT om verschillende banen in één assenstelsel te
tekenen.
\end{question}

\end{oefening}


% ============================================================
% 11 PROBLEEMOPLOSSEN
% ============================================================

\FVDsection{Een horizontale worp mathematiseren}

Bij een horizontale worp is het zelden verstandig meteen naar één
formule te zoeken.

Werk systematisch:

\begin{enumerate}
  \item Maak een schets.
  \item Kies een oorsprong en positieve richtingen.
  \item Splits de beweging op in $x$ en $y$.
  \item Noteer voor elke richting de beginpositie, beginsnelheid en
        versnelling.
  \item Noteer wat op het gezochte tijdstip bekend is.
  \item Gebruik de componentbeweging waarin voldoende informatie
        beschikbaar is om eerst $t$ te bepalen.
  \item Gebruik diezelfde tijd in de andere componentbeweging.
  \item Stel indien nodig de snelheidsvector samen.
  \item Controleer teken, eenheid, grootteorde en modelaannames.
\end{enumerate}

\begin{eigenschap}

Dezelfde tijd $t$ verschijnt in beide componentbewegingen. Dat is vaak
de sleutel om een horizontale-worpprobleem op te lossen.

\end{eigenschap}


\begin{oefening}[Reddingspakket]{\oefU\ESS}

Een drone vliegt horizontaal met constante snelheid

\[
v_0=18,0\,\mathrm{m/s}
\]

op een hoogte van $45,0\,\mathrm{m}$. Een klein pakket wordt
losgelaten. Verwaarloos luchtweerstand.

\begin{question}
Hoe lang valt het pakket?
\end{question}

\begin{question}
Hoe ver vóór het doel moet het pakket worden losgelaten?
\end{question}

\begin{question}
Bepaal de horizontale en verticale snelheidscomponent bij impact.
\end{question}

\begin{question}
Bepaal de grootte van de impactsnelheid.
\end{question}

\begin{question}
Leg uit waarom het pakket bij het loslaten een horizontale
beginsnelheid heeft, hoewel het alleen maar wordt ``losgelaten''.
\end{question}

\end{oefening}


\begin{oefening}[Stuntman]{\oefG}

Een stuntman rijdt horizontaal van een platform en moet op een lager
platform landen.

Het landingsplatform ligt $3,20\,\mathrm{m}$ lager en de horizontale
afstand tussen de randen bedraagt $8,50\,\mathrm{m}$.

\begin{question}
Bepaal de minimale horizontale snelheid waarmee hij de rand van het
landingsplatform bereikt volgens het puntmassamodel.
\end{question}

\begin{question}
Bepaal de grootte en richting van zijn snelheid op dat ogenblik.
\end{question}

\begin{question}
Welke aspecten van een echte stunt worden door dit model genegeerd?
Noem er minstens drie.
\end{question}

\end{oefening}


% ============================================================
% 12 LABO
% ============================================================

\FVDsection{Labo: onderzoek een horizontale worp}

Een horizontale worp leent zich uitstekend voor videoanalyse.

\textbf{Onderzoek experimenteel het verband tussen positie, tijd, snelheid en versnelling bij een horizontale worp.}



Een mogelijke werkwijze:

\begin{enumerate}
  \item Laat een kleine bal horizontaal van een tafel rollen.
  \item Film de beweging loodrecht op het bewegingsvlak.
  \item Zorg voor een lengteschaal in beeld.
  \item Kies in de videoanalyse een oorsprong en assenstelsel.
  \item Registreer op opeenvolgende tijdstippen $x$ en $y$.
  \item Verwerk de meetgegevens grafisch.
\end{enumerate}


\FVDsubsection{Horizontale analyse}

Maak een grafiek van $x$ als functie van $t$.

Volgens het model verwachten we:

\[
x(t)=x_0+v_0t.
\]

Dus een lineaire trendlijn.

De helling levert een schatting van $v_x$.


\FVDsubsection{Verticale analyse}

Maak een grafiek van $y$ als functie van $t$.

Volgens het model verwachten we:

\[
y(t)=y_0-\frac12gt^2.
\]

Dus een kwadratisch verband.

Bij een regressiemodel

\[
y(t)=At^2+Bt+C
\]

verwachten we idealiter:

\[
A\approx-\frac g2,
\qquad
B\approx0,
\qquad
C\approx y_0.
\]


\begin{oefening}[Laboanalyse]{\oefU\ESS}

Verwerk de meetgegevens van de horizontale worp.

\begin{question}
Maak een $x(t)$-grafiek en bepaal een geschikte trendlijn.
\end{question}

\begin{question}
Interpreteer de parameters fysisch.
\end{question}

\begin{question}
Maak een $y(t)$-grafiek en bepaal een kwadratische trendlijn.
\end{question}

\begin{question}
Bepaal uit het kwadratische model een experimentele waarde voor $g$.
\end{question}

\begin{question}
Maak indien mogelijk ook $v_x(t)$ en $v_y(t)$ zichtbaar.
\end{question}

\begin{question}
Vergelijk experiment en theorie en bespreek mogelijke
afwijkingsbronnen.
\end{question}

\end{oefening}


% ============================================================
% 13 MODELKRITIEK
% ============================================================

\FVDsection{Het ideale projectielmodel}

Onze horizontale worp is een model.

We nemen onder andere aan:

\begin{itemize}
  \item het voorwerp is een puntmassa;
  \item luchtweerstand is verwaarloosbaar;
  \item $g$ is constant;
  \item de aarde kan lokaal als vlak worden beschouwd;
  \item rotatie en aerodynamische effecten zijn niet relevant.
\end{itemize}

Voor een kleine compacte bal over een beperkte afstand kan dat een
zeer bruikbaar model zijn. Voor een pluimbal, een voetbal met sterke
spin of een projectiel over zeer grote afstand kunnen de afwijkingen
belangrijk worden.


\begin{oefening}[Model of werkelijkheid?]{\oefU}

Vergelijk de geschiktheid van het ideale horizontale-worpmodel voor:

\begin{question}
een stalen knikker die van een tafel rolt;
\end{question}

\begin{question}
een prop papier die horizontaal wordt weggegooid;
\end{question}

\begin{question}
een voetbal met sterke spin;
\end{question}

\begin{question}
een waterdruppel uit een tuinslang.
\end{question}

Geef telkens aan welke modelaanname mogelijk problematisch is.

\end{oefening}


% ============================================================
% 14 GEINTEGREERDE KERN
% ============================================================

\FVDsection{Geïntegreerde oefeningen horizontale worp}

\begin{oefening}[Van functie naar fysica]{\oefU\ESS}

Van een horizontale worp worden de componentfuncties gegeven:

\[
x(t)=2,0+9,0t,
\]

\[
y(t)=15,0-4,905t^2.
\]

Alle afstanden zijn in meter en de tijd in seconde.

\begin{question}
Bepaal de beginpositie.
\end{question}

\begin{question}
Bepaal de beginsnelheidsvector.
\end{question}

\begin{question}
Bepaal de versnellingsvector.
\end{question}

\begin{question}
Bepaal wanneer het voorwerp $y=0$ bereikt.
\end{question}

\begin{question}
Bepaal de horizontale positie op dat ogenblik.
\end{question}

\begin{question}
Bepaal de snelheidsvector en de grootte van de snelheid bij impact.
\end{question}

\end{oefening}


\begin{oefening}[Ontwerp achterwaarts]{\oefG}

Een horizontaal gelanceerde bal moet vanaf een platform van
$5,00\,\mathrm{m}$ hoogte een doel raken dat zich
$7,50\,\mathrm{m}$ horizontaal van de platformrand bevindt.

\begin{question}
Bepaal de vereiste horizontale beginsnelheid.
\end{question}

\begin{question}
Stel de baanvergelijking op.
\end{question}

\begin{question}
Bepaal de snelheid waarmee de bal het doel bereikt.
\end{question}

\begin{question}
Als de lanceersnelheid een relatieve onzekerheid van $2\%$ heeft,
bespreek kwalitatief hoe dit de horizontale landingspositie beïnvloedt.
\end{question}

\end{oefening}


\begin{oefening}[Redeneer zonder rekenen]{\oefU\ESS}

Twee ballen worden horizontaal van dezelfde hoogte gelanceerd.
Bal B heeft een driemaal zo grote horizontale beginsnelheid als bal A.

Verwaarloos luchtweerstand.

Vergelijk zonder eerst getallen te kiezen:

\begin{question}
de valtijden;
\end{question}

\begin{question}
de verticale impactsnelheden;
\end{question}

\begin{question}
de horizontale drachten;
\end{question}

\begin{question}
de groottes van de totale impactsnelheden.
\end{question}

Motiveer elke conclusie.

\end{oefening}


% ============================================================
% 15 VERDIEPING SCHUINE WORP
% ============================================================

\FVDsection{Verdieping: de schuine worp}

\begin{opmerking}

Deze sectie is een \textbf{uitbreiding}. De horizontale worp behoort
tot de kern van dit hoofdstuk. De schuine worp gaat verder en is
bedoeld voor leerlingen die de kern vlot beheersen en extra uitdaging
willen.

\end{opmerking}

Bij een schuine worp maakt de beginsnelheid een hoek $\alpha$ met de
horizontale.

We ontbinden:

\[
\boxed{
v_{0x}=v_0\cos\alpha
}
\]

en:

\[
\boxed{
v_{0y}=v_0\sin\alpha.
}
\]

Daarna verandert het principe niet.

Horizontaal:

\[
a_x=0
\]

en:

\[
\boxed{
x(t)=x_0+v_0\cos\alpha\,t.
}
\]

Verticaal:

\[
a_y=-g
\]

en:

\[
\boxed{
y(t)
=
y_0+v_0\sin\alpha\,t-\frac12gt^2.
}
\]

De snelheidscomponenten zijn:

\[
\boxed{
v_x(t)=v_0\cos\alpha
}
\]

en:

\[
\boxed{
v_y(t)=v_0\sin\alpha-gt.
}
\]

De horizontale worp is dus het bijzondere geval:

\[
\boxed{\alpha=0^\circ.}
\]


\FVDsubsection{Het hoogste punt}

Op het hoogste punt:

\[
v_y=0.
\]

Dus:

\[
v_0\sin\alpha-gt_{\max}=0
\]

en:

\[
\boxed{
t_{\max}
=
\frac{v_0\sin\alpha}{g}.
}
\]

De horizontale snelheidscomponent is daar niet nul:

\[
v_x=v_0\cos\alpha.
\]

De totale snelheid is op het hoogste punt dus horizontaal.


\FVDsubsection{Maximale hoogtewinst}

Voor de verticale component kunnen we het resultaat uit de verticale
worp gebruiken:

\[
\boxed{
\Delta y_{\max}
=
\frac{v_0^2\sin^2\alpha}{2g}.
}
\]


\FVDsubsection{Baanvergelijking}

Uit:

\[
x-x_0=v_0\cos\alpha\,t
\]

volgt:

\[
t
=
\frac{x-x_0}{v_0\cos\alpha}.
\]

Invullen in $y(t)$ geeft:

\[
\boxed{
y(x)
=
y_0
+
(x-x_0)\tan\alpha
-
\frac{g(x-x_0)^2}
{2v_0^2\cos^2\alpha}.
}
\]

Ook de ideale schuine worp volgt dus een parabool.


\begin{oefening}[Een schuine worp analyseren]{\oefU}

Een bal wordt vanaf de grond gelanceerd met:

\[
v_0=20,0\,\mathrm{m/s},
\qquad
\alpha=35^\circ.
\]

Verwaarloos luchtweerstand.

\begin{question}
Bepaal $v_{0x}$ en $v_{0y}$.
\end{question}

\begin{question}
Stel $x(t)$ en $y(t)$ op.
\end{question}

\begin{question}
Bepaal het tijdstip en de hoogte van het hoogste punt.
\end{question}

\begin{question}
Bepaal de snelheidsvector op het hoogste punt.
\end{question}

\end{oefening}


% ============================================================
% 16 VERDIEPING GELIJKE HOOGTE
% ============================================================

\FVDsection{Verdieping: gelijke begin- en eindhoogte}

Wanneer een schuine worp eindigt op dezelfde hoogte als waarop hij
begon, kunnen enkele elegante relaties worden afgeleid.

Naast de triviale oplossing $t=0$ geldt voor de terugkeer:

\[
0
=
v_0\sin\alpha\,t-\frac12gt^2.
\]

Dus:

\[
\boxed{
T
=
\frac{2v_0\sin\alpha}{g}.
}
\]

De dracht is:

\[
R
=
v_0\cos\alpha\,T.
\]

Invullen geeft:

\[
R
=
\frac{2v_0^2\sin\alpha\cos\alpha}{g}.
\]

Met

\[
2\sin\alpha\cos\alpha=\sin(2\alpha)
\]

volgt:

\[
\boxed{
R
=
\frac{v_0^2\sin(2\alpha)}{g}.
}
\]


\FVDsubsection{Welke hoek geeft de grootste dracht?}

Voor vaste $v_0$ en gelijke begin- en eindhoogte is $R$ maximaal
wanneer:

\[
\sin(2\alpha)=1.
\]

Dus:

\[
2\alpha=90^\circ
\]

en:

\[
\boxed{
\alpha=45^\circ.
}
\]

Dit resultaat geldt binnen het ideale model. In echte sporten of
technische toepassingen kunnen luchtweerstand, rotatie, lanceerhoogte
en aerodynamica de optimale hoek veranderen.


\begin{oefening}[Complementaire hoeken]{\oefG}

Toon aan dat twee lanceerhoeken

\[
\alpha
\qquad\text{en}\qquad
90^\circ-\alpha
\]

bij dezelfde $v_0$ en gelijke begin- en eindhoogte dezelfde
horizontale dracht geven.

Gebruik zowel de formule voor $R$ als een inhoudelijke redenering over
de horizontale en verticale component van de beginsnelheid.

\end{oefening}


\begin{oefening}[Is $45^\circ$ altijd optimaal?]{\oefG}

Een leerling beweert:

\begin{quote}
``Een projectiel vliegt altijd het verst onder een hoek van
$45^\circ$.''
\end{quote}

Analyseer deze uitspraak.

\begin{question}
Onder welke aannames is ze correct?
\end{question}

\begin{question}
Waarom hoeft ze niet te gelden wanneer begin- en eindhoogte
verschillen?
\end{question}

\begin{question}
Waarom hoeft ze niet te gelden in een echte sportcontext?
\end{question}

\end{oefening}


% ============================================================
% 17 VERDIEPING MODELLEREN
% ============================================================

\FVDsection{Verdieping: ontwerp een projectielbaan}

\begin{oefening}[Van doel naar lanceerhoek]{\oefG}

Een projectiel wordt vanaf de oorsprong gelanceerd met een vaste
beginsnelheid $v_0$. Een doel bevindt zich in het punt $(X,Y)$.

Uit de baanvergelijking:

\[
Y
=
X\tan\alpha
-
\frac{gX^2}{2v_0^2\cos^2\alpha}
\]

kan je onderzoeken welke lanceerhoeken het doel kunnen bereiken.

\begin{question}
Leg uit waarom er voor sommige doelen twee mogelijke lanceerhoeken
kunnen bestaan.
\end{question}

\begin{question}
Leg uit waarom sommige doelen met de gegeven $v_0$ helemaal niet
bereikbaar zijn.
\end{question}

\begin{question}
Gebruik ICT om voor een zelfgekozen $v_0$ een familie van banen voor
verschillende $\alpha$ te tekenen.
\end{question}

\begin{question}
Formuleer op basis van de grafieken een hypothese over het bereikbare
gebied en onderzoek die verder.
\end{question}

\end{oefening}


% ============================================================
% 18 SAMENVATTING
% ============================================================

\FVDsection{Wat hebben we opgebouwd?}

Bij een tweedimensionale beweging:

\[
\boxed{
\vec r=x\,\vec e_x+y\,\vec e_y
}
\]

\[
\boxed{
\vec v=v_x\,\vec e_x+v_y\,\vec e_y
}
\]

\[
\boxed{
\vec a=a_x\,\vec e_x+a_y\,\vec e_y.
}
\]

De componentbewegingen kunnen afzonderlijk worden geanalyseerd, maar
delen dezelfde tijd.

Voor een horizontale worp met omhoog positief:

\[
\boxed{
a_x=0,
\qquad
a_y=-g
}
\]

\[
\boxed{
x(t)=x_0+v_0t
}
\]

\[
\boxed{
y(t)=y_0-\frac12gt^2
}
\]

\[
\boxed{
v_x=v_0,
\qquad
v_y=-gt.
}
\]

De totale snelheid:

\[
\boxed{
|\vec v|
=
\sqrt{v_x^2+v_y^2}.
}
\]

Vanaf hoogte $h$:

\[
\boxed{
t_{\mathrm{val}}
=
\sqrt{\frac{2h}{g}}
}
\]

en:

\[
\boxed{
d
=
v_0\sqrt{\frac{2h}{g}}.
}
\]

De baanvergelijking:

\[
\boxed{
y(x)
=
y_0-\frac{g}{2v_0^2}(x-x_0)^2.
}
\]

De ideale baan is dus parabolisch.


% ============================================================
% 20 VOORUITBLIK
% ============================================================

\FVDsection{Vooruitblik}

Bij de horizontale worp veranderde de snelheidsvector zowel in grootte
als in richting.

Er bestaat ook een beweging waarbij de grootte van de snelheid
constant blijft, maar de richting voortdurend verandert: de
eenparige cirkelbeweging.

Ook daar zal de vectoriële beschrijving essentieel blijken. De
versnellingsvector wijst dan niet in de bewegingsrichting, maar naar
het middelpunt van de cirkel.

\end{document}
